% Options for packages loaded elsewhere
\PassOptionsToPackage{unicode}{hyperref}
\PassOptionsToPackage{hyphens}{url}
\documentclass[
]{article}
\usepackage{xcolor}
\usepackage[margin=1in]{geometry}
\usepackage{amsmath,amssymb}
\setcounter{secnumdepth}{-\maxdimen} % remove section numbering
\usepackage{iftex}
\ifPDFTeX
  \usepackage[T1]{fontenc}
  \usepackage[utf8]{inputenc}
  \usepackage{textcomp} % provide euro and other symbols
\else % if luatex or xetex
  \usepackage{unicode-math} % this also loads fontspec
  \defaultfontfeatures{Scale=MatchLowercase}
  \defaultfontfeatures[\rmfamily]{Ligatures=TeX,Scale=1}
\fi
\usepackage{lmodern}
\ifPDFTeX\else
  % xetex/luatex font selection
\fi
% Use upquote if available, for straight quotes in verbatim environments
\IfFileExists{upquote.sty}{\usepackage{upquote}}{}
\IfFileExists{microtype.sty}{% use microtype if available
  \usepackage[]{microtype}
  \UseMicrotypeSet[protrusion]{basicmath} % disable protrusion for tt fonts
}{}
\makeatletter
\@ifundefined{KOMAClassName}{% if non-KOMA class
  \IfFileExists{parskip.sty}{%
    \usepackage{parskip}
  }{% else
    \setlength{\parindent}{0pt}
    \setlength{\parskip}{6pt plus 2pt minus 1pt}}
}{% if KOMA class
  \KOMAoptions{parskip=half}}
\makeatother
\usepackage{graphicx}
\makeatletter
\newsavebox\pandoc@box
\newcommand*\pandocbounded[1]{% scales image to fit in text height/width
  \sbox\pandoc@box{#1}%
  \Gscale@div\@tempa{\textheight}{\dimexpr\ht\pandoc@box+\dp\pandoc@box\relax}%
  \Gscale@div\@tempb{\linewidth}{\wd\pandoc@box}%
  \ifdim\@tempb\p@<\@tempa\p@\let\@tempa\@tempb\fi% select the smaller of both
  \ifdim\@tempa\p@<\p@\scalebox{\@tempa}{\usebox\pandoc@box}%
  \else\usebox{\pandoc@box}%
  \fi%
}
% Set default figure placement to htbp
\def\fps@figure{htbp}
\makeatother
% definitions for citeproc citations
\NewDocumentCommand\citeproctext{}{}
\NewDocumentCommand\citeproc{mm}{%
  \begingroup\def\citeproctext{#2}\cite{#1}\endgroup}
\makeatletter
 % allow citations to break across lines
 \let\@cite@ofmt\@firstofone
 % avoid brackets around text for \cite:
 \def\@biblabel#1{}
 \def\@cite#1#2{{#1\if@tempswa , #2\fi}}
\makeatother
\newlength{\cslhangindent}
\setlength{\cslhangindent}{1.5em}
\newlength{\csllabelwidth}
\setlength{\csllabelwidth}{3em}
\newenvironment{CSLReferences}[2] % #1 hanging-indent, #2 entry-spacing
 {\begin{list}{}{%
  \setlength{\itemindent}{0pt}
  \setlength{\leftmargin}{0pt}
  \setlength{\parsep}{0pt}
  % turn on hanging indent if param 1 is 1
  \ifodd #1
   \setlength{\leftmargin}{\cslhangindent}
   \setlength{\itemindent}{-1\cslhangindent}
  \fi
  % set entry spacing
  \setlength{\itemsep}{#2\baselineskip}}}
 {\end{list}}
\usepackage{calc}
\newcommand{\CSLBlock}[1]{\hfill\break\parbox[t]{\linewidth}{\strut\ignorespaces#1\strut}}
\newcommand{\CSLLeftMargin}[1]{\parbox[t]{\csllabelwidth}{\strut#1\strut}}
\newcommand{\CSLRightInline}[1]{\parbox[t]{\linewidth - \csllabelwidth}{\strut#1\strut}}
\newcommand{\CSLIndent}[1]{\hspace{\cslhangindent}#1}
\setlength{\emergencystretch}{3em} % prevent overfull lines
\providecommand{\tightlist}{%
  \setlength{\itemsep}{0pt}\setlength{\parskip}{0pt}}
\usepackage{bookmark}
\IfFileExists{xurl.sty}{\usepackage{xurl}}{} % add URL line breaks if available
\urlstyle{same}
\hypersetup{
  pdftitle={peppwR: Simulation-based power analysis for phosphoproteomics experiments},
  pdfauthor={Dan MacLean, The Sainsbury Laboratory, Norwich Research Park, Norwich NR4 7UH, UK},
  hidelinks,
  pdfcreator={LaTeX via pandoc}}

\title{peppwR: Simulation-based power analysis for phosphoproteomics
experiments}
\author{Dan MacLean, The Sainsbury Laboratory, Norwich Research Park,
Norwich NR4 7UH, UK}
\date{}

\begin{document}
\maketitle
\begin{abstract}
\textbf{Background:} Phosphoproteomics experiments require careful
experimental design to ensure adequate statistical power, yet no
dedicated tools exist for power analysis that account for the unique
statistical properties of peptide-level quantification data.

\textbf{Results:} We present peppwR, an R package for simulation-based
power analysis tailored to phosphoproteomics. peppwR fits distributions
to pilot data on a per-peptide basis, capturing the heterogeneity in
abundance and variance across the peptidome. It supports multiple
statistical tests (Wilcoxon, bootstrap-t, Bayes factor), tracks
missingness patterns common in mass spectrometry data, and provides
FDR-aware power calculations. Users can determine required sample sizes,
estimate power for planned experiments, or identify minimum detectable
effect sizes.

\textbf{Availability:} peppwR is freely available from GitHub
(\url{https://github.com/TeamMacLean/peppwR}) with documentation at
\url{https://teammaclean.github.io/peppwR/}.
\end{abstract}

\begin{figure}
\centering
\includegraphics[width=1\linewidth,height=\textheight,keepaspectratio]{figure_1.pdf}
\caption{\textbf{Figure 1. peppwR workflow and example results.} (A)
Two-stage workflow: fit distributions to pilot data, then run power
simulations. (B) Distribution of per-peptide power at N=6 samples per
group for detecting 2-fold changes. The dashed line indicates 80\% power
threshold. Power varies dramatically across peptides due to differences
in abundance, variance, and missingness. (C) Power curves comparing
Wilcoxon and Bayes factor tests. Y-axis shows the proportion of peptides
that individually achieve 80\% power at each sample size. Bayes factor
reaches the 80\% threshold (dotted line) at N=10 versus N=20 for
Wilcoxon. (D) Power lookup table across sample sizes and effect sizes.
Values indicate the percentage of peptides achieving 80\% power for each
combination.}
\end{figure}

\section{Introduction}\label{introduction}

Statistical power---the probability of detecting a true effect when one
exists---is fundamental to experimental design. Under-powered
experiments waste resources and risk missing genuine biological signals,
while over-powered experiments are inefficient (Cohen 1988; Button et
al. 2013). Power analysis is standard practice in clinical trials and
genomics, yet remains neglected in proteomics despite the field's
increasing quantitative sophistication.

Generic power analysis tools fail to capture the unique statistical
properties of mass spectrometry-based proteomics data. First, peptides
exhibit substantial heterogeneity: abundance levels span orders of
magnitude, and variance differs dramatically across the peptidome (Olsen
et al. 2006; Humphrey, Azimifar, and Mann 2015). Second, abundance
distributions are typically right-skewed rather than normal, better
described by gamma or lognormal distributions. Third, missing values
follow non-random patterns---low-abundance peptides are systematically
more likely to be undetected, creating missing-not-at-random (MNAR) data
(Webb-Robertson et al. 2015; Lazar et al. 2016). Fourth,
phosphoproteomics experiments simultaneously test thousands of peptides,
requiring consideration of multiple testing correction.

Existing tools address these challenges incompletely. Perseus (Tyanova
et al. 2016), the most widely-used proteomics analysis platform,
provides comprehensive statistical analysis but lacks power analysis
functionality. The clippda package (Nyangoma et al. 2012) was designed
for SELDI-TOF peak data and assumes homogeneous variance across
features. MultiPower (Tarazona et al. 2020) focuses on multi-omics
integration, requiring multiple data types. G*Power (Faul et al. 2007)
is a general-purpose tool that assumes normality and lacks
proteomics-specific features.

We present peppwR, an R package that addresses these limitations through
per-peptide distribution fitting, simulation-based power estimation, and
missingness-aware calculations. peppwR answers three key experimental
design questions: ``What sample size do I need?'', ``What power will I
achieve?'', and ``What is the minimum detectable effect size?''

\section{Implementation}\label{implementation}

\subsection{Workflow}\label{workflow}

peppwR implements a two-stage workflow (Figure 1A). First,
\texttt{fit\_distributions()} models the abundance distribution for each
peptide in pilot data, selecting the best-fitting distribution from
gamma, lognormal, normal, and inverse Gaussian candidates using AIC.
Second, \texttt{power\_analysis()} runs Monte Carlo simulations drawing
from these fitted distributions to estimate power.

Two operational modes accommodate different use cases. In
\textbf{aggregate mode}, users specify a single distribution and
parameters, enabling power calculations without pilot data. In
\textbf{per-peptide mode}, peppwR leverages pilot data to model
heterogeneity across the peptidome, reporting the proportion of peptides
achieving target power at each sample size.

\subsection{Distribution fitting}\label{distribution-fitting}

For each peptide, peppwR fits candidate distributions using maximum
likelihood estimation via the fitdistrplus (Delignette-Muller and Dutang
2015) and univariateML packages. The best-fitting distribution is
selected by minimum AIC. When fitting fails (e.g., insufficient
non-missing observations), users can choose to exclude the peptide, fall
back to empirical bootstrap resampling, or use a moment-matched
lognormal distribution.

\subsection{Simulation engine}\label{simulation-engine}

Power is estimated via Monte Carlo simulation. For each peptide and each
candidate sample size:

\begin{enumerate}
\def\labelenumi{\arabic{enumi}.}
\tightlist
\item
  Draw \emph{n} samples from the fitted control distribution
\item
  Draw \emph{n} samples from a treatment distribution scaled by the
  specified effect size
\item
  Apply the selected statistical test
\item
  Record whether the null hypothesis is rejected
\item
  Repeat for \emph{n\textsubscript{sim}} iterations
\end{enumerate}

Power is the proportion of iterations achieving statistical
significance. In per-peptide mode, peppwR reports the proportion of
peptides exceeding target power at each sample size.

\subsection{Statistical tests}\label{statistical-tests}

peppwR supports three statistical tests. The \textbf{Wilcoxon rank-sum
test} (Wilcoxon 1945) (default) is non-parametric and robust to
distributional assumptions. The \textbf{bootstrap-t test} (Efron and
Tibshirani 1994) handles non-normality through resampling. The
\textbf{Bayes factor t-test} (Rouder et al. 2009) provides a Bayesian
approach that quantifies evidence strength rather than dichotomous
rejection, offering improved efficiency when parametric assumptions
hold.

\subsection{Missing data handling}\label{missing-data-handling}

peppwR tracks missingness at two levels. Per-peptide, it records the
proportion of missing values. At the dataset level, it detects MNAR
(missing-not-at-random) patterns by correlating mean abundance with
missingness rate across peptides---a negative correlation indicates that
low-abundance peptides are systematically more likely to be missing, the
hallmark of detection-limit-driven MNAR common in mass spectrometry
(Webb-Robertson et al. 2015). When
\texttt{include\_missingness\ =\ TRUE}, simulations incorporate observed
per-peptide NA rates, providing power estimates that account for
expected data loss.

\subsection{FDR-aware power}\label{fdr-aware-power}

For whole-peptidome studies, peppwR can simulate complete experiments
with Benjamini-Hochberg FDR correction (Benjamini and Hochberg 1995).
Setting \texttt{apply\_fdr\ =\ TRUE} runs simulations across all
peptides simultaneously, reporting the proportion of true positives
discovered at the target FDR threshold.

\section{Example Application}\label{example-application}

We demonstrate peppwR using simulated phosphoproteomics data mimicking
typical pilot study characteristics: 500 peptides with heterogeneous
gamma-distributed abundances and \textasciitilde13\% missingness with a
realistic MNAR pattern (see Supplementary Material for complete code).

\subsection{Per-peptide power
analysis}\label{per-peptide-power-analysis}

Using \texttt{fit\_distributions()} followed by
\texttt{power\_analysis()} with \texttt{find\ =\ "sample\_size"}, we
determined sample size requirements for detecting 2-fold changes using
the Wilcoxon test. Because each peptide has its own distribution
parameters, power is calculated separately for each peptide via
simulation. We then ask: at what sample size do 80\% of peptides
individually achieve 80\% power? This threshold-based approach
acknowledges that not all peptides will be equally detectable. The
answer: N = 20 samples per group.

\subsection{Power heterogeneity}\label{power-heterogeneity}

Figure 1B illustrates the central insight motivating per-peptide
analysis: at N = 6 samples per group, power varies dramatically across
peptides, ranging from near-zero to over 90\%. This heterogeneity arises
from differences in abundance, variance, and missingness. A single
``average'' power estimate would obscure this crucial variability and
potentially mislead experimental design decisions.

\subsection{Effect of statistical test
choice}\label{effect-of-statistical-test-choice}

Comparing Wilcoxon and Bayes factor tests reveals substantial
differences in sample size requirements (Figure 1C). The Bayes factor
t-test achieves the 80\% threshold at N = 10, compared to N = 20 for
Wilcoxon---a 50\% reduction in required samples. This efficiency gain
reflects the Bayes factor's parametric assumptions, which are reasonable
when the underlying distributions are well-characterized by pilot data.
Users should match their power analysis to their planned statistical
approach.

\subsection{Power lookup table}\label{power-lookup-table}

Figure 1D presents a practical lookup table showing the proportion of
peptides achieving 80\% power across combinations of sample size and
effect size. At N = 12 with 3-fold changes, 94\% of peptides achieve
adequate power; at N = 4 with 1.5-fold changes, only 1\% do. Such tables
enable researchers to make informed trade-offs between detectable effect
sizes and experimental costs.

\section{Discussion}\label{discussion}

peppwR provides the first dedicated power analysis tool for
phosphoproteomics that accounts for per-peptide heterogeneity,
non-normal distributions, and missingness patterns. The simulation-based
approach accommodates the complex statistical properties of mass
spectrometry data that violate assumptions of traditional power
calculations.

Our example analysis demonstrates that power varies substantially across
the peptidome, and that analytical choices---particularly the
statistical test---meaningfully impact sample size requirements. By
providing per-peptide power estimates, peppwR enables researchers to set
realistic expectations about which peptides will be detectable at their
planned sample size.

\textbf{Limitations.} Computational cost scales linearly with peptide
count and number of simulations; typical analyses complete within
minutes on modern hardware. Per-peptide mode requires pilot data from
the same experimental system; when unavailable, aggregate mode with
reasonable distributional assumptions provides useful guidance.

\textbf{Future directions.} Integration with the pepdiff package for
differential abundance analysis would provide an end-to-end workflow
from experimental design through analysis. Extension to TMT and iTRAQ
labeling schemes, which have different variance structures, represents a
natural next step.

\section{Funding}\label{funding}

This work was supported by the Gatsby Charitable Foundation.

\section{Conflict of Interest}\label{conflict-of-interest}

None declared.

\section*{References}\label{references}
\addcontentsline{toc}{section}{References}

\protect\phantomsection\label{refs}
\begin{CSLReferences}{1}{0}
\bibitem[\citeproctext]{ref-benjamini1995fdr}
Benjamini, Yoav, and Yosef Hochberg. 1995. {``Controlling the False
Discovery Rate: A Practical and Powerful Approach to Multiple
Testing.''} \emph{Journal of the Royal Statistical Society: Series B
(Methodological)} 57 (1): 289--300.
\url{https://doi.org/10.1111/j.2517-6161.1995.tb02031.x}.

\bibitem[\citeproctext]{ref-button2013power}
Button, Katherine S, John PA Ioannidis, Claire Mokrysz, Brian A Nosek,
Jonathan Flint, Emma SJ Robinson, and Marcus R Munafò. 2013. {``Power
Failure: Why Small Sample Size Undermines the Reliability of
Neuroscience.''} \emph{Nature Reviews Neuroscience} 14 (5): 365--76.
\url{https://doi.org/10.1038/nrn3475}.

\bibitem[\citeproctext]{ref-cohen1988power}
Cohen, Jacob. 1988. {``Statistical Power Analysis for the Behavioral
Sciences.''} \url{https://doi.org/10.4324/9780203771587}.

\bibitem[\citeproctext]{ref-delignette2015fitdistrplus}
Delignette-Muller, Marie Laure, and Christophe Dutang. 2015.
{``{fitdistrplus}: An {R} Package for Fitting Distributions.''}
\emph{Journal of Statistical Software} 64 (4): 1--34.
\url{https://doi.org/10.18637/jss.v064.i04}.

\bibitem[\citeproctext]{ref-efron1994bootstrap}
Efron, Bradley, and Robert J. Tibshirani. 1994. \emph{An Introduction to
the Bootstrap}. 1st ed. Chapman; Hall/CRC.
\url{https://doi.org/10.1201/9780429246593}.

\bibitem[\citeproctext]{ref-gpower}
Faul, Franz, Edgar Erdfelder, Albert-Georg Lang, and Axel Buchner. 2007.
{``{G*Power} 3: A Flexible Statistical Power Analysis Program for the
Social, Behavioral, and Biomedical Sciences.''} \emph{Behavior Research
Methods} 39 (2): 175--91. \url{https://doi.org/10.3758/bf03193146}.

\bibitem[\citeproctext]{ref-humphrey2015phosphoproteomics}
Humphrey, Sean J, S Barzin Azimifar, and Matthias Mann. 2015.
{``High-Throughput Phosphoproteomics Reveals in Vivo Insulin Signaling
Dynamics.''} \emph{Nature Biotechnology} 33 (9): 990--95.
\url{https://doi.org/10.1038/nbt.3327}.

\bibitem[\citeproctext]{ref-lazar2016mnar}
Lazar, Cosmin, Laurent Gatto, Myriam Ferro, Christophe Bruley, and
Thomas Burger. 2016. {``Accounting for the Multiple Natures of Missing
Values in Label-Free Quantitative Proteomics Data Sets to Compare
Imputation Strategies.''} \emph{Journal of Proteome Research} 15 (4):
1116--25. \url{https://doi.org/10.1021/acs.jproteome.5b00981}.

\bibitem[\citeproctext]{ref-nyangoma2012clippda}
Nyangoma, Stephen O, Stuart I Collins, Douglas G Altman, Philip Johnson,
and Lucinda Billingham. 2012. {``Sample Size Calculations for Designing
Clinical Proteomic Profiling Studies Using Mass Spectrometry.''}
\emph{Statistical Applications in Genetics and Molecular Biology} 11
(3). \url{https://doi.org/10.1515/1544-6115.1686}.

\bibitem[\citeproctext]{ref-olsen2006phosphoproteomics}
Olsen, Jesper V, Blagoy Blagoev, Florian Gnad, Boris Macek, Chanchal
Kumar, Peter Mortensen, and Matthias Mann. 2006. {``Global, in Vivo, and
Site-Specific Phosphorylation Dynamics in Signaling Networks.''}
\emph{Cell} 127 (3): 635--48.
\url{https://doi.org/10.1016/j.cell.2006.09.026}.

\bibitem[\citeproctext]{ref-rouder2009bayest}
Rouder, Jeffrey N, Paul L Speckman, Dongchu Sun, Richard D Morey, and
Geoffrey Iverson. 2009. {``Bayesian t Tests for Accepting and Rejecting
the Null Hypothesis.''} \emph{Psychonomic Bulletin \& Review} 16 (2):
225--37. \url{https://doi.org/10.3758/PBR.16.2.225}.

\bibitem[\citeproctext]{ref-tarazona2020multipower}
Tarazona, Sonia, Leandro Balzano-Nogueira, David Gómez-Cabrero, Andreas
Schmidt, Sara C Madeira, and Ana Conesa. 2020. {``Harmonization of
Quality Metrics and Power Calculation in Multi-Omic Studies.''}
\emph{Nature Communications} 11 (1): 3092.
\url{https://doi.org/10.1038/s41467-020-16937-8}.

\bibitem[\citeproctext]{ref-tyanova2016perseus}
Tyanova, Stefka, Tikira Temu, Pavel Sinitcyn, Arthur Carlson, Marco Y
Hein, Tamar Geiger, Matthias Mann, and Jürgen Cox. 2016. {``The
{Perseus} Computational Platform for Comprehensive Analysis of
(Prote)omics Data.''} \emph{Nature Methods} 13 (9): 731--40.
\url{https://doi.org/10.1038/nmeth.3901}.

\bibitem[\citeproctext]{ref-webb2015mnar}
Webb-Robertson, Bobbie-Jo M, Holli K Wiber, Melissa M Matzke, Jemima N
Samuel, Karin D Rodland, Thomas RW Clauss, and Joel G Pounds. 2015.
{``Review, Evaluation, and Discussion of the Challenges of Missing Value
Imputation for Mass Spectrometry-Based Label-Free Global Proteomics.''}
\emph{Journal of Proteome Research} 14 (5): 1993--2001.
\url{https://doi.org/10.1021/pr501138h}.

\bibitem[\citeproctext]{ref-wilcoxon1945}
Wilcoxon, Frank. 1945. {``Individual Comparisons by Ranking Methods.''}
\emph{Biometrics Bulletin} 1 (6): 80--83.
\url{https://doi.org/10.2307/3001968}.

\end{CSLReferences}

\end{document}
